\begin{comment}
keep toc to d=1 for this section --- reduces ToC page size. default is d=3
\end{comment}

\addtocontents{toc}{\protect\setcounter{tocdepth}{2}}

\subsection{Buterin's Trilemma}
\label{subsec:trilemma}

\href{\citeShardingFAQsLink}{The trilemma} is as follows:

Given $c$ (computational resources per node, e.g., computation, bandwidth, storage, etc.), and $n$ (the size of the network, e.g., transaction throughput, state size, user population, market cap, etc.), blockchain systems have no more than two of these three properties:

\begin{itemize}
  \item Decentralization --- the system can operate with participants that have only $O(c)$ resources (e.g., a laptop, a raspberry pi, a VPS; typical internet bandwidth, storage space, etc.).
  \item Security --- the system is secure against attackers with up to $O(n)$ resources.\footnotemark
  \item Scalability --- the system can process $O(n)$ transactions, with $O(n) > O(c)$; this means that, as the network grows, the throughput of the system grows faster than the computational resources required per user.
\end{itemize}

\footnotetext{
  Here, by convention, ``up to $O(n)$ resources'' means that the system remains secure against attackers with less than some critical proportion of the network's block production capacity --- e.g., less than $\nicefrac{1}{2}$ for typical PoW chains, and less than $\nicefrac{1}{3}$ for typical PoS chains.
}

This definition of scalability is perhaps problematic.
If the network, which has $O(n)$ TPS demand, becomes bottlenecked by an $O(c^2)$ scaling method, is the network really scalable?
I prefer an alternative definition of scalability: the system can process $O(n)$ transactions in $O(1)$ time, i.e., getting a transaction confirmed is neither too expensive nor delayed as $n$ and/or $c$ change.\footnote{
    This does not mean there is no fee market for transactions, nor that localized congestion will not happen.
    Rather, we are interested in an equilibrium that provides substantial value and capacity whilst avoiding transactions being as free as sending an email.
}

Why mention $O(c^2)$ scaling particularly?
Why is that an important breakpoint for scaling configurations?
In a word: sharding.
The standard method of sharding (or hosting child-chains generally) is to replace transactions with shard-headers in the host-chain.
Extra data might also be required.
If the host-chain has $O(c)$ capacity, then it should support\footnote{
    Presuming a secure method of sharding is known and in use, and $O(1)$ load per shard-header.
} $O(c)$ shards.
Each shard has $O(c)$ capacity also, thus the full system has $O(c^2)$ capacity.

\aside{
  Traditionally, the use of big $O$ notation when discussing the trilemma has not conformed to the strict mathematical meaning; instead, it is used quite loosely as a way to analyze the scalability of a blockchain architecture.
  Naturally, we'll continue this tradition.
}

\subsubsection{Core Conflict}

\label{sec:core-conflict}

\begin{figure}
\centering
\includegraphics[max width=\linewidth]{trilemma/core_conflict_sag}
\captionsetup{justification=centering,singlelinecheck=false}
\caption{
    A ``cloud'' diagram of the core conflict of \textit{Buterin's Trilemma}.\\
    $A \to B$ can be read ``$A$ so that $B$''; and $A \leftarrow B$: ``$A$ requires $B$''.
}
\label{fig:trilemma-core-conflict}
\end{figure}

\autoref{fig:trilemma-core-conflict} shows a cloud\footnote{
    Clouds are diagrams of explanatory implications which reveal a conflict between two (apparently) necessary things.
    For more on clouds, see \href{https://www.google.com/search?q=ISBN+0-88427-115-3}{It's Not Luck} by Eli Goldratt (1994).
} for the core conflict.
It reads: \emph{safely increasing capacity} requires that we \emph{stay decentralized}.
\emph{Staying decentralized} requires that we \emph{use many chains and small blocks}.
\emph{Safely increasing capacity} requires that the network \emph{stays secure}.
\emph{Staying secure} requires that we \emph{use big blocks}.
Big blocks are the opposite of small blocks, so we have a \emph{conflict}.
Additionally: using multiple chains compromises security; and big blocks compromise decentralization.

To understand the core conflict, we need to look at the underlying assumptions.

Buterin writes\footnote{See \href{\citeShardingFAQsLink}{\emph{On sharding blockchains FAQs}}.} in the Ethereum sharding FAQ regarding the \emph{many chains with small blocks} strategy:

\begin{quote}
{[}\ldots{}{]} This greatly increases throughput, but at a cost of
security: an N-factor increase in throughput using this method
necessarily comes with an N-factor decrease in security, as a level of
resources 1/N the size of the whole ecosystem will be sufficient to
attack any individual chain. {[}\ldots{}{]}
\end{quote}

He writes, regarding the \emph{big blocks} strategy:

\begin{quote}
{[}\ldots{}{]} such an approach inevitably has its limits: if one goes
too far, then nodes running on consumer hardware will drop out, the
network will start to rely exclusively on a very small number of
supercomputers running the blockchain, which can lead to great
centralization risk.
\end{quote}

A mistaken way to break the conflict is \emph{merged mining} (a.k.a.
AuxPoW). This method attempts to share security between chains, so that
one might be able to have decentralization and security via small blocks
and merged mined chains/shards (sometimes called side-chains). Buterin
writes:

\begin{quote}
{[}\ldots{}{]} If all miners participate, this theoretically can
increase throughput by a factor of N without compromising security.
However, this also has the problem that it increases the computational
and storage load on each miner by a factor of N, and so in fact this
solution is simply a stealthy form of block size increase.
\end{quote}

\autoref{fig:trilemma-mm-conflict} shows the cloud for the merged mining
conflict.


\begin{figure}[H]
\centering
\includegraphics[max width=\linewidth]{trilemma/mm_conflict_sag}
\caption{A cloud showing the scaling conflict of \textit{merged mining}.}
\label{fig:trilemma-mm-conflict}
\end{figure}

An underlying assumption here is that maximally sharing security across the network requires miners to maintain a record of all chains and do validation on all those chains.
The naive solution (using many chains, mentioned above) conflicts with secure methods of merged mining.
It seems like progress is impossible.

We have some hints to conditions that might belong to a solution:

\begin{itemize}
  \item We can share security between chains/shards and can use small blocks.
  \item We can share security without miners keeping and validating all chains/shards.
\end{itemize}

\textbf{This is the crux of the problem: how do you construct a network of blockchains (scalable) such that attacking an individual chain is about as difficult as attacking the full network (secure), whilst ensuring that the security of the network does not require validating all chains (decentralized)?}

\subsubparagraph{Prior Assumptions}

Here are some prior underlying assumptions that are either common or which I expect to be:

\begin{itemize}
  \item Sharing PoW security requires merged mining.
  \item Sharing PoW security requires that chains use the same hashing algorithm.
  \item Multiple (non-merged-mined) PoW chains using the same hashing algorithm means that at least some of those chains are vulnerable to a 51\% attack.
  \item Simultaneously securing a network with PoW and PoS is not possible without compromises (like that PoW miners could DoS PoS validators or vice versa).
  \item It is unsafe for miners/validators to build on unvalidated histories (as is done with SPV mining, which \emph{is} unsafe).
\end{itemize}

\begin{comment}
These assumptions \emph{are} true in lots of cases (arguably in all cases up to now).
Are they always?
\end{comment}

\aside{
    I call them \emph{assumptions}; some people will likely (and rightly) take issue with that and call them \emph{conclusions} instead.
    For our purposes there isn't really a difference; I include them here so that I can later show you conditions under which they are all simultaneously false.
}

\subsubsection{Conjecture: A Principle of Scaling}

\label{sec:a-principle-of-scaling}

A scalable system \emph{can} have components with complexities worse than $O(c)$ \emph{if and only if} those components are not \emph{bottlenecks}, i.e., \emph{constraints}.
As long as there is \emph{excess capacity} in those components, the system can still scale.
%%\text{(See also: \emph{Theory of Constraints}; \emph{Critical Fallibilism}.)}

\begin{comment}
end of section, reset toc to default
\end{comment}

\SetToCDepthDefault
